% ============================================================
% TEACHING STATEMENT
% Standalone document. Compile with pdfLaTeX.
% ============================================================

\documentclass[11pt,letterpaper]{article}

\usepackage[
    margin=0.75in,
    headheight=14pt,
    headsep=14pt,
    footskip=22pt
]{geometry}

\usepackage[T1]{fontenc}
\usepackage[utf8]{inputenc}
\usepackage{newtxtext}
\usepackage{microtype}
\usepackage{xcolor}
\usepackage{titlesec}
\usepackage{fancyhdr}
\usepackage[hidelinks]{hyperref}

% ============================================================
% COLORS
% ============================================================

\definecolor{Accent}{HTML}{007F6B}
\definecolor{Body}{HTML}{222222}
\definecolor{Muted}{HTML}{555555}

% ============================================================
% SPACING AND TYPOGRAPHY
% ============================================================

\setlength{\parindent}{0pt}
\setlength{\parskip}{3.5pt}
\linespread{1.0}

\raggedbottom
\emergencystretch=1.5em
\clubpenalty=10000
\widowpenalty=10000

\titleformat{\section}
    {\large\bfseries\color{Accent}}
    {}
    {0pt}
    {}

\titlespacing*{\section}
    {0pt}
    {8pt}
    {3pt}

\titleformat{\paragraph}[runin]
    {\normalsize\bfseries}
    {}
    {0pt}
    {}
    [.]

\titlespacing*{\paragraph}
    {0pt}
    {5pt}
    {0.5em}

% ============================================================
% HEADERS AND FOOTERS
% ============================================================

\pagestyle{fancy}
\fancyhf{}

\fancyhead[L]
    {\small\color{Muted}Mohammad Alipour-Vaezi}

\fancyhead[R]
    {\small\color{Muted}Teaching Statement}

\fancyfoot[C]
    {\small\color{Muted}\thepage}

\renewcommand{\headrulewidth}{0.3pt}
\renewcommand{\footrulewidth}{0pt}

\fancypagestyle{firstpage}{%
    \fancyhf{}
    \fancyfoot[C]{\small\color{Muted}\thepage}
    \renewcommand{\headrulewidth}{0pt}
}

\hypersetup{
    pdftitle={Teaching Statement},
    pdfauthor={Mohammad Alipour-Vaezi}
}

% ============================================================
% DOCUMENT
% ============================================================

\begin{document}

\color{Body}
\thispagestyle{firstpage}

{\LARGE\bfseries Teaching Statement}
\par
\vspace{3pt}

{\large Mohammad Alipour-Vaezi}
\par

{\small\color{Muted}
Industrial and Systems Engineering \textbar{} Virginia Tech}
\par
\vspace{5pt}

{\color{Accent}\hrule height 0.6pt}
\vspace{6pt}

I want students to turn important practical questions into defensible
decisions. That requires more than solving a model or obtaining an accurate
prediction: students must identify the decision, explain their assumptions,
evaluate alternatives, and communicate the limits of their evidence. My
teaching connects operations research, business analytics, and artificial
intelligence. My central goal is quantitative independence: the ability to
choose and evaluate a method rather than rely on a familiar formula or
software output.

\section{Teaching experience and preparation}

I have served as an instructor for \textbf{Deterministic Operations Research}
at Virginia Tech in Summer 2026 and at Shahid Beheshti University of Medical
Sciences in Summer 2022, and for \textbf{Introduction to Industrial and
Systems Engineering} at Virginia Tech in Fall 2024.

My graduate teaching assistant experience broadens that foundation. At
Virginia Tech, my assignments have included Deterministic Operations
Research, Optimization I, Simulation I, Data Management for Industrial and
Systems Engineering, Introduction to Industrial and Systems Engineering,
and Manufacturing Process Lab. At the University of Tehran, I supported
Scheduling Theory, Queuing Systems, Design of Industrial Systems, and
Advanced Engineering Economics.

My preparation includes Virginia Tech's \textbf{Graduate Certificate of
Future Professoriate}. In 2025, I received the \textbf{White GTA recognition}
from the Grado Department of Industrial and Systems Engineering and the
\textbf{Outstanding GTA award} from the Alpha Pi Mu Honor Society. I will
build on this preparation by assessing what students can explain,
implement, and transfer to unfamiliar problems.

\section{Connecting modeling, computation, and judgment}

\paragraph{Begin with the decision, then develop the method}
I would organize a resource-allocation case around its verbal description,
mathematical formulation, and computational implementation. Students would
identify variables and units, formulate operational constraints, solve a
small instance, and implement a larger model. They would then change a
capacity or cost assumption and explain the resulting recommendation. In a
business classroom, the deliverable would include a managerial decision
memo; in an advanced engineering classroom, it would also include a
derivation or optimality argument.

\paragraph{Make reasoning visible}
I would combine short explanations with worked examples, individual
attempts, and discussion of competing formulations. Comparing two plausible
models can reveal assumptions that are otherwise easy to overlook. A
simulation assignment would similarly require justification of the
experimental design and interpretation of uncertainty, not merely a
numerical estimate. Accessible starting problems and analytical extensions
would allow students with different preparation to work toward the same
conceptual objective.

\paragraph{Assess understanding as well as execution}
My assessment plan combines individual problem solving, computational
assignments, and projects with staged feedback. Projects would progress
from a problem statement to a baseline analysis, a revised model, and a
final recommendation. Rubrics would distinguish formulation,
implementation, validation, and interpretation. Individual explanations
of team projects would make each student's understanding visible.
Low-stakes checks and midsemester feedback would identify topics requiring
another explanation and inform improvements to subsequent offerings.

\section{Courses I am prepared to teach}

My strongest immediate contributions are \textbf{operations research and
optimization, simulation, and data-driven decision-making}. At the
undergraduate level, I am prepared to teach Introduction to Industrial and
Systems Engineering, Deterministic Operations Research, Linear and Integer
Optimization, Simulation, and introductory Business Analytics or Management
Science. My teaching experience also supports contributions to Data
Management, Scheduling, Queuing Systems, and Engineering Economics.

At the graduate level, my research prepares me to teach Reinforcement
Learning, Sequential Decision-Making, Risk-Sensitive Optimization, and
advanced prescriptive analytics. For professional business students, I
would emphasize formulation, interpretation, and managerial tradeoffs;
for engineering and doctoral students, mathematical analysis, algorithm
design, and reproducible implementation would receive greater weight.

\section{Proposed courses}

I would develop the following electives in response to departmental needs,
initially as special-topics offerings that complement the existing
curriculum.

\paragraph{Reinforcement Learning for Operations and Decision-Making}
This graduate elective would connect Markov decision processes, dynamic
programming, exploration, and policy evaluation to inventory, service
operations, and sequential resource allocation. Students would examine
how policies change when reliability matters alongside average reward.
A final project would require reproducible implementation, comparison
with a simple operational policy, and an explanation of the assumptions
supporting the result. Prerequisites would include probability,
optimization, and introductory programming.

\paragraph{Causal and Prescriptive Analytics for Healthcare}
This advanced undergraduate or master's course would distinguish predicting
outcomes from estimating intervention effects and choosing actions. Topics
would include study design, confounding, heterogeneous responses, resource
constraints, and multi-attribute decisions. Using public or synthetic cases,
students would analyze patient prioritization, treatment comparisons, or
capacity allocation. Their final deliverable would combine a decision memo,
an analytical notebook, and an explicit account of uncertainty. The
emphasis would be responsible analysis, not clinical recommendations.

\paragraph{Reliable AI Agents for Business Operations}
This advanced undergraduate or master's elective would examine retrieval,
memory, tool use, and human oversight in language-model systems. Students
would evaluate an agent as an operational system with costs, failure
modes, and escalation requirements. Projects would compare agents with
simpler alternatives under changing information and matched computational
budgets. Assessment would emphasize reproducible testing and justified
limits on automation. Public or synthetic tasks would avoid dependence on
proprietary systems.

\section{Access, responsible tool use, and mentorship}

Students with different mathematical and programming backgrounds need clear
routes into quantitative work. I would provide prerequisite refreshers,
annotated examples, accessible materials, and transparent assessment
criteria. Written questions and small-group discussions would complement
whole-class participation. Office hours would support both students
encountering difficulty and those pursuing deeper extensions. Feedback
would identify the next step in a student's reasoning rather than simply
label an answer correct or incorrect.

My generative artificial intelligence policy would depend on the learning
objective. Foundational assessments would establish independent
understanding; selected projects would allow disclosed assistance with
coding or exploration. Students would remain responsible for checking
calculations, verifying references, protecting restricted data, and
explaining their submissions. An unexplained correct output would not
substitute for demonstrated understanding.

My leadership as president and vice president of the Virginia Tech INFORMS
Student Chapter provides a foundation for contributing beyond the
classroom. As a faculty mentor, I would use regular meetings, written
milestones, reproducible research practices, and early discussions of
expectations and authorship. Undergraduate projects could begin with
replication or a focused modeling question; doctoral mentoring would
progressively transfer responsibility for problem selection, analysis,
and scholarly communication. I want students to become colleagues who can
formulate worthwhile questions and defend their answers with clarity
and care.

\end{document}